\pagestyle{frontmatter}
\pagenumbering{roman}

% -------------------- DECLARATION --------------------
\chapter*{DECLARATION}
\addcontentsline{toc}{chapter}{Declaration}

I, \textbf{Jiphin George}, student of Master of Computer Applications, Department of Computer Applications, Mar Athanasius College of Engineering, Kothamangalam, hereby declare that this Internship Report titled \textbf{``QualiVision AI – AI Based Live Tomato Freshness Detection System''} is a bonafide record of the internship work carried out by me at \textbf{Nestsoft Technomaster Pvt. Ltd.}, during the period \textbf{25 May 2026} to \textbf{03 July 2026}, in partial fulfilment of the requirements for the award of the degree of Master of Computer Applications.

I further declare that this report has not been submitted previously, either in part or in full, to any other University or Institute for the award of any degree or diploma.

\vspace{1.5cm}

\noindent Place: Kothamangalam\\
\noindent Date: \textbf{03 July 2026} \hfill \textbf{Jiphin George}

\newpage

% -------------------- ACKNOWLEDGEMENT --------------------
\chapter*{ACKNOWLEDGEMENT}
\addcontentsline{toc}{chapter}{Acknowledgement}

I express my sincere gratitude to the Almighty for providing me with the strength, wisdom, and guidance to complete this academic internship project and this report successfully.

I extend my heartfelt thanks to Prof. Sonia Abraham, Head of the Department of Computer Applications, Mar Athanasius College of Engineering, Kothamangalam, for the constant encouragement, institutional support, and motivation throughout my academic programme.

I am deeply grateful to Prof. Shinu S Kurian, my Internship Guide, and Prof. Sibu Skaria, my Internship Coordinator, in the Department of Computer Applications. I sincerely thank them for their valuable guidance, timely feedback, mentorship, and for efficiently facilitating the internship programme and review processes.

I sincerely thank the mentors and team at Nestsoft Technomaster Pvt. Ltd. for the technical guidance, professional support, constructive suggestions, and the opportunity to work on the QualiVision AI project.

I express my gratitude to all the faculty members and non-teaching staff of the Department of Computer Applications for their continued support and cooperation throughout my studies.

Finally, I am deeply thankful to my parents, family, and friends for their unwavering encouragement, patience, and moral support at every stage of this journey.

\vspace{1cm}
\begin{flushright}
\textbf{Jiphin George}
\end{flushright}

\newpage

% -------------------- ABSTRACT --------------------
\chapter*{ABSTRACT}
\addcontentsline{toc}{chapter}{Abstract}

This report presents a comprehensive overview of the six-week internship undertaken at Nestsoft Technomaster Pvt. Ltd. as part of the Master of Computer Applications (MCA) programme at Mar Athanasius College of Engineering, Kothamangalam, from 25 May 2026 to 03 July 2026. The internship offered an extensive, hands-on learning journey through Artificial Intelligence, Data Science, Machine Learning, and Computer Vision, bridging the gap between theoretical academic concepts and professional software engineering.

The internship commenced with foundational concepts in Artificial Intelligence, Python programming, and the development of basic applications like an AI Translator and a Text-to-Speech system. It progressively transitioned into comprehensive Data Science practices, focusing on Exploratory Data Analysis (EDA) and visualization using Pandas, Matplotlib, and Seaborn. Advanced machine learning modules introduced supervised and unsupervised learning algorithms, allowing for predictive modeling on varied datasets, including student performance and passenger survival predictions.

The culmination of this internship was the capstone project: ``QualiVision AI – AI Based Live Tomato Freshness Detection System.'' Designed to automate quality inspection in food processing industries, the system leverages a convolutional neural network (EfficientNetV2B0) to classify agricultural produce as fresh or rotten. It features real-time prediction capabilities, a robust Flask backend, an SQLite database for persistent logging, and a professional industrial dashboard created using Google Stitch. Furthermore, Gradient-weighted Class Activation Mapping (Grad-CAM) was implemented to provide Explainable AI (XAI) insights, ensuring transparency in the model's decision-making process.

Through the successful completion of these modules and the QualiVision AI project, the internship significantly enhanced practical competencies in AI model development, deployment, backend engineering, and industrial automation, laying a solid foundation for a successful career in Artificial Intelligence and Software Development.

\newpage
